%% kapa: exam=1 round=24 year=2013 period=2 form=A subject=accounting subject_ko=회계학 seq=35 number=35
%% source: Q-Net 감정평가사 기출문제 (cst00309) artlSeq=5073802 file=extracted/1008/5073802/2013년 제24회 감정평가사 1차시험 기출문제/2교시(A형).hwp
%% conversion: kordoc 4.24.51 -> tools/gen1.py (auto); review=proofread-sonnet; flags=none
\begin{question}{35}
(주)한국은 동일한 공정에서 A, B, C라는 3가지 결합제품을 생산하고 있다. 결합원가 ￦1,600은 제품별 순실현가치에 비례하여 배부한다. 제품별 자료가 다음과 같을 때, 제품 단위당 매출총이익이 높은 것부터 낮은 순으로 나열한 것은? \par
{\small\centering
\begin{tabular}{|c|c|c|c|}
\hline
\makecell{제품명} & \makecell{A} & \makecell{B} & \makecell{C} \\
\hline
\makecell{생산량 및 판매량} & \makecell{100단위} & \makecell{300단위} & \makecell{500단위} \\
\hline
\makecell{단위당 판매가격} & \makecell{￦13} & \makecell{￦10} & \makecell{￦15} \\
\hline
\makecell{단위당 추가가공원가} & \makecell{3} & \makecell{2} & \makecell{-} \\
\hline
\makecell{단위당 판매비} & \makecell{1} & \makecell{1} & \makecell{5} \\
\hline
\end{tabular}\par}

\begin{choices}
  \item A > B > C
  \item A > C > B
  \item C > B > A
  \item C > A > B
  \item B > C > A
\end{choices}
\end{question}
\answer{4}
